% =====================================================================
% Module 2: To chuc luu tru du lieu toan he thong
% ---------------------------------------------------------------------
% Cach su dung:
%   1. Dat file .tex nay cung thu muc voi bao cao chinh, hoac \input{} no.
%   2. Mo file "module2_diagrams.html" bang trinh duyet, chup man hinh tung
%      bieu do (Hinh 2.1 -> 2.2) va luoi thanh anh (PNG/PDF) dat cung thu muc
%      voi ten: module2_hinh2_1.png ... module2_hinh2_2.png
%   3. Chup man hinh S3 console that cua ban va luu thanh:
%      module2_hinh2_3.png  (Bronze bucket - danh sach source folders)
%      module2_hinh2_4.png  (Bronze - drill-down year/month/day + .jsonl.gz)
%      module2_hinh2_5.png  (Silver bucket - danh sach source_site partitions)
%      module2_hinh2_6.png  (Silver - drill-down + .parquet files)
%      module2_hinh2_7.png  (MotherDuck UI - Gold database tables)
%      module2_hinh2_8.png  (Qdrant Cloud UI - collection jobs_v1)
%   4. Go bo khoi "lstlisting setup" ben duoi neu bao cao da co package
%      "listings" va dinh nghia kieu Python.
% =====================================================================

% ---- Vietnamese encoding (xoa neu da co trong preamble) -------------
\usepackage[utf8]{inputenc}
\usepackage[T5]{fontenc}
% ---- lstlisting setup (xoa neu da co trong preamble) ----------------
\usepackage{listings}
\usepackage{xcolor}
\lstdefinestyle{pystyle}{
    language=Python,
    basicstyle=\ttfamily\footnotesize,
    keywordstyle=\color{blue}\bfseries,
    stringstyle=\color{teal},
    commentstyle=\color{gray}\itshape,
    showstringspaces=false,
    breaklines=true,
    frame=single,
    rulecolor=\color{lightgray},
    backgroundcolor=\color{white},
    tabsize=4,
    columns=fullflexible,
    keepspaces=true,
}
\lstdefinestyle{yamlstyle}{
    language={},
    basicstyle=\ttfamily\footnotesize,
    keywordstyle=\color{blue}\bfseries,
    stringstyle=\color{teal},
    commentstyle=\color{gray}\itshape,
    showstringspaces=false,
    breaklines=true,
    frame=single,
    rulecolor=\color{lightgray},
    backgroundcolor=\color{white},
    tabsize=2,
    columns=fullflexible,
    keepspaces=true,
}
\lstset{style=pystyle}
% ---------------------------------------------------------------------

\usepackage{graphicx}
\usepackage{booktabs}
\usepackage{longtable}
\usepackage{hyperref}


\section{Module 2: Tổ chức lưu trữ dữ liệu toàn hệ thống}
\label{sec:storage-all}

Sau khi Crawl Layer ghi dữ liệu thô ra file JSONL tạm cục bộ (Module~1), Storage
Layer tiếp quản việc nén, upload lên AWS S3, làm sạch, và phân phối dữ liệu đến
các hệ thống downstream: MotherDuck (Gold, BI), Supabase (serving frontend), và
Qdrant (vector search cho gợi ý việc làm theo CV). Hệ thống áp dụng kiến trúc
\textbf{Medallion} (Bronze -- Silver -- Gold) trên hai S3 bucket riêng biệt, mỗi
lớp có định dạng file và quy ước đường dẫn khác nhau. Bên cạnh Medallion trên
S3, Qdrant lưu vector embedding từ Silver để phục vụ semantic search.

\subsection{Tổng quan kiến trúc Medallion}
\label{subsec:medallion-overview}

Kiến trúc Medallion chia dữ liệu thành ba lớp chất lượng tăng dần:

\begin{itemize}
    \item \textbf{Bronze} --- dữ liệu thô y nguyên từ crawler, nén gzip ở định
    dạng \texttt{.jsonl.gz}. Lớp này chỉ đảm bảo \emph{tồn tại} và \emph{không
    sửa đổi} dữ liệu gốc.
    \item \textbf{Silver} --- dữ liệu đã làm sạch, chuẩn hoá, lưu ở định dạng
    Parquet (\texttt{.parquet}) cho hiệu suất truy vấn cột tốt hơn. Đây là lớp
    trung tâm mà mọi downstream đều đọc từ đó.
    \item \textbf{Gold} --- star schema trên MotherDuck (DuckDB cloud), gồm fact
    table \texttt{gold.jobs} và các dimension/child table. Phục vụ BI (Power BI).
\end{itemize}

Bên cạnh Gold, Silver còn feeding trực tiếp vào \textbf{Supabase} (Postgres) ---
lớp serving cho frontend thông qua UPSERT trên \texttt{job\_url}, và vào
\textbf{Qdrant} --- vector database cho gợi ý việc làm theo CV. Sơ đồ tổng thể
được thể hiện ở Hình~\ref{fig:medallion-arch}.

\begin{figure}[htbp]
    \centering
    % Chup man hinh tu module2_diagrams.html -> Hinh 2.1
    \includegraphics[width=\textwidth]{module2_hinh2_1.png}
    \caption{Kiến trúc lưu trữ toàn hệ thống --- dòng chảy từ Crawl Layer
    qua hai bucket Bronze/Silver, đến Gold (MotherDuck), Supabase, và Qdrant
    (vector DB). Phần trong khung nét đứt là không gian AWS S3.}
    \label{fig:medallion-arch}
\end{figure}


\subsection{Cấu hình bucket S3}
\label{subsec:bucket-config}

Hệ thống sử dụng \textbf{hai S3 bucket} riêng biệt, tên bucket được hardcode
trong file YAML duy nhất để dễ đổi khi đổi AWS account:

\begin{lstlisting}[style=yamlstyle, caption={\texttt{src/storage\_layer/MinIO\_S3/config/bucket.yml}}, label={lst:bucket-yml}]
bucket_name:
  bronze_layer: thethien-lakehouse-lite-bronze
  silver_layer: thethien-lakehouse-lite-silver
\end{lstlisting}

Tên bucket phải \emph{globally unique} trên toàn AWS (không chỉ trong account),
nên được prefix bằng tên dự án \texttt{thethien-lakehouse-lite-}. Khi chia sẻ
AWS account, cần đổi tên bucket trong file này.

Hình~\ref{fig:s3-bronze-bucket} chụp màn hình S3 Console thực tế hiển thị Bronze bucket và danh sách
các thư mục source (itviec, topcv, vietnamworks).

\begin{figure}[htbp]
    \centering
    % CHUP MAN HINH S3 CONSOLE: Bucket thethien-lakehouse-lite-bronze
    % Hien thi danh sach folder: itviec/, topcv/, vietnamworks/
    \includegraphics[width=\textwidth]{module2_hinh2_3.png}
    \caption{S3 Console --- Bronze bucket \texttt{thethien-lakehouse-lite-bronze}
    với các thư mục source (itviec, topcv, vietnamworks).}
    \label{fig:s3-bronze-bucket}
\end{figure}


\subsection{Kết nối S3 qua boto3}
\label{subsec:s3-connection}

Điểm truy cập S3 duy nhất trong toàn codebase là hàm \texttt{get\_s3\_client()}
trong \texttt{minio\_connect.py}. Dù tên file vẫn giữ \texttt{minio\_connect} (di
sản từ khi dự án dùng MinIO), hàm này trả về \textbf{boto3 client thật} kết nối
AWS S3, không có \texttt{endpoint\_url} override:

\begin{lstlisting}[caption={\texttt{src/storage\_layer/MinIO\_S3/utils/minio\_connect.py}}, label={lst:minio-connect}]
import boto3
from src.storage_layer.MinIO_S3.config.key import (
    AWS_ACCESS_KEY_ID, AWS_SECRET_ACCESS_KEY, AWS_REGION,
)

# Module name kept as `minio_connect` to avoid touching ~30 import sites.
# This now returns a real AWS S3 client; no endpoint_url override.
def get_s3_client():
    return boto3.client(
        "s3",
        aws_access_key_id=AWS_ACCESS_KEY_ID,
        aws_secret_access_key=AWS_SECRET_ACCESS_KEY,
        region_name=AWS_REGION,
    )
\end{lstlisting}

Credential được đọc từ \texttt{.env} tại import time trong \texttt{key.py}:

\begin{lstlisting}[caption={\texttt{src/storage\_layer/MinIO\_S3/config/key.py}}, label={lst:key-py}]
from dotenv import load_dotenv
import os

load_dotenv()

AWS_ACCESS_KEY_ID = os.getenv("AWS_ACCESS_KEY_ID")
AWS_SECRET_ACCESS_KEY = os.getenv("AWS_SECRET_ACCESS_KEY")
AWS_REGION = os.getenv("AWS_REGION", "ap-southeast-1")
\end{lstlisting}

Mọi module cần thao tác S3 (Bronze upload, Silver reader, Silver loader,
MotherDuck) đều import \texttt{get\_s3\_client()} --- đây là \emph{single S3
client factory} của toàn hệ thống.


\subsection{Quy ước đường dẫn (Hive partitioning)}
\label{subsec:hive-partitioning}

S3 không có thư mục thật --- ``thư mục'' là ảo, được suy ra từ dấu \texttt{/}
trong object key. Hệ thống dùng \textbf{Hive-style partitioning}: các segment
có dạng \texttt{key=value} (vd: \texttt{year=2026/month=06/day=21}) để
DuckDB/Polars tự nhận diện partition columns khi scan.

Hai lớp Bronze và Silver có \textbf{quy ước khác nhau có chủ đích}:

\begin{itemize}
    \item \textbf{Bronze} dùng prefix \texttt{source/} --- thư mục đơn giản,
    phù hợp vì Bronze chỉ lưu raw, không cần partition pruning.
    \item \textbf{Silver} dùng \texttt{source\_site=} --- Hive partition thật,
    cho phép DuckDB \texttt{read\_parquet(hive\_partitioning=true)} tự expose
    \texttt{source\_site} thành cột, hỗ trợ filter theo site mà không đọc file.
\end{itemize}

Logic sinh đường dẫn tập trung trong hai class \texttt{BronzeBucketPaths} và
\texttt{SilverBucketPaths}:

\begin{lstlisting}[caption={\texttt{src/storage\_layer/MinIO\_S3/config/path.py} --- sinh S3 prefix cho Bronze và Silver}, label={lst:path-py}]
DEFAULT_ENTITY_NAME = "jobs"

class BronzeBucketPaths:
    def __init__(self, source_name, entity_name=DEFAULT_ENTITY_NAME,
                 year="*", month="*", day="*"):
        # ... load bucket name from bucket.yml ...
        self.bronze_bucket_name = self._get_bronze_bucket_name()

    def get_prefix(self) -> str:
        return f"{self.source_name}/{self.entity_name}/" \
               f"year={self.year}/month={self.month}/day={self.day}/"

    def get_files_path_json_gz(self) -> str:
        return f"s3://{self.bronze_bucket_name}/{self.get_prefix()}*.jsonl.gz"


class SilverBucketPaths:
    def __init__(self, source_site, entity_name=DEFAULT_ENTITY_NAME,
                 year="*", month="*", day="*"):
        self.silver_bucket_name = self._get_silver_bucket_name()

    def get_prefix(self) -> str:
        return f"{self.entity_name}/source_site={self.source_site}/" \
               f"year={self.year}/month={self.month}/day={self.day}/"

    def get_file_key(self, timestamp: str) -> str:
        return f"{self.get_prefix()}clean_bronze_{timestamp}.parquet"
\end{lstlisting}

Hình~\ref{fig:hive-compare} so sánh trực quan cấu trúc đường dẫn của hai lớp.

\begin{figure}[htbp]
    \centering
    % Chup man hinh tu module2_diagrams.html -> Hinh 2.2
    \includegraphics[width=\textwidth]{module2_hinh2_2.png}
    \caption{So sánh cấu trúc Hive partitioning: Bronze dùng \texttt{source/}
    prefix, Silver dùng \texttt{source\_site=} Hive key. Sự khác biệt này cho
    phép DuckDB nhận diện \texttt{source\_site} thành partition column khi scan
    Silver parquet.}
    \label{fig:hive-compare}
\end{figure}


\subsection{Bronze Layer: Nén và upload}
\label{subsec:bronze-upload}

Hàm \texttt{load\_to\_bronze()} trong \texttt{loader.py} (thuộc Crawl Layer
utils) thực hiện toàn bộ luồng Bronze: quét temp dir, nén gzip, upload S3,
xoá file tạm. Đây là thao tác \textbf{destructive} --- sau khi upload thành
công, cả file \texttt{.gz} lẫn \texttt{.jsonl} gốc đều bị xoá:

\begin{lstlisting}[caption={\texttt{src/crawl\_layer/utils/loader.py} --- \texttt{load\_to\_bronze()} (trích đoạn chính)}, label={lst:load-to-bronze}]
def load_to_bronze(bucket_name="bronze", source_filter=None):
    s3_client = get_s3_client()
    # Create bucket if it doesn't exist
    try:
        s3_client.head_bucket(Bucket=bucket_name)
    except Exception:
        s3_client.create_bucket(
            Bucket=bucket_name,
            CreateBucketConfiguration={
                'LocationConstraint': s3_client.meta.region_name
            }
        )

    timestamp = datetime.now().strftime("%H%M%S")

    for file_path in TEMP_DIR.glob("*.jsonl"):
        # Parse: itviec_jobs_20260512.jsonl -> source, entity, date
        parts = file_path.stem.split('_')
        source_name, entity_name, date_str = parts[0], parts[1], parts[2]
        year, month, day = date_str[:4], date_str[4:6], date_str[6:8]

        # Compress .jsonl -> .jsonl.gz
        gz_file_path = file_path.with_suffix('.jsonl.gz')
        with open(file_path, 'rb') as f_in:
            with gzip.open(gz_file_path, 'wb') as f_out:
                f_out.writelines(f_in)

        # S3 key: source/entity/year=Y/month=M/day=D/name_HHMMSS.jsonl.gz
        s3_file_name = f"{source_name}_{entity_name}_{date_str}_{timestamp}.jsonl.gz"
        s3_key = f"{source_name}/{entity_name}/year={year}/month={month}/day={day}/{s3_file_name}"
        s3_client.upload_file(str(gz_file_path), bucket_name, s3_key)

        gz_file_path.unlink()  # xoa file .gz sau khi upload

    # Xoa tat ca .jsonl cua source (hoac toan bo neu khong co filter)
    clean_temp_directory(prefix=source_filter)
\end{lstlisting}

Các điểm quan trọng:
\begin{itemize}
    \item Timestamp \texttt{\%H\%M\%S} được thêm vào tên file để \emph{tránh
    ghi đè} khi chạy lại cùng ngày --- mỗi lần upload tạo object mới.
    \item Tham số \texttt{source\_filter} cho phép xử lý \emph{một site} cụ
    thể, bỏ qua file của site khác. Khi \texttt{None}, xử lý toàn bộ.
    \item \texttt{clean\_temp\_directory(prefix=source\_filter)} chỉ xoá file
    khớp prefix, để các site khác không bị ảnh hưởng.
\end{itemize}

Hình~\ref{fig:s3-bronze-drill} chụp màn hình S3 Console cho thấy cấu trúc partition
\texttt{year=2026/month=06/day=21/} thật và các file \texttt{.jsonl.gz} bên trong.

\begin{figure}[htbp]
    \centering
    % CHUP MAN HINH S3 CONSOLE: Bronze bucket drill-down
    % Vao itviec/jobs/year=2026/month=06/day=21/
    % Hien thi cac file .jsonl.gz
    \includegraphics[width=\textwidth]{module2_hinh2_4.png}
    \caption{S3 Console --- Bronze bucket drill-down theo Hive partition
    \texttt{year=2026/month=06/day=21/}, hiển thị các file
    \texttt{.jsonl.gz} đã nén. Mỗi file là một lần upload (timestamp khác nhau).}
    \label{fig:s3-bronze-drill}
\end{figure}


\subsection{Silver Layer: Làm sạch và Parquet}
\label{subsec:silver-layer}

Silver Layer đọc Bronze JSONL, làm sạch, và ghi ra Parquet trên S3 bucket thứ
hai. Đây là lớp phức tạp nhất với bốn thành phần chính: data model, reader,
loader, và cleaning pipeline.

\subsubsection{Data model: SilverJobItem}
\label{subsubsec:silver-data-model}

Schema Silver được định nghĩa duy nhất trong dataclass \texttt{SilverJobItem}
--- \emph{single source of truth}. Thêm/xoá field ở đây, schema Polars tự sinh:

\begin{lstlisting}[caption={\texttt{src/storage\_layer/MinIO\_S3/layer/silver/data\_model/data\_class.py} --- trích đoạn}, label={lst:silver-job-item}]
@dataclass
class SilverJobItem:
    # 1. METADATA & SOURCE
    job_url: Optional[str] = None
    unique_url: Optional[str] = None
    search_keyword: Optional[str] = None
    job_deadline: Optional[str] = None

    # 2. JOB TITLE
    job_title: Optional[str] = None
    clean_job_title: Optional[str] = None
    job_title_special_keywords: List[str] = field(default_factory=list)

    # 3. COMPANY INFO
    clean_company_name: Optional[str] = None
    company_size: Optional[str] = None
    min_company_size: Optional[int] = None
    max_company_size: Optional[int] = None

    # 4-10. LOCATION, INDUSTRY, SALARY, DESCRIPTION, REQUIREMENTS ...
    # (30+ fields tong cong, chia 10 nhom)

# Auto-generate Polars schema from dataclass
def silver_schema_to_polars() -> dict[str, pl.DataType]:
    hints = get_type_hints(SilverJobItem)
    schema: dict[str, pl.DataType] = {}
    for name, tp in hints.items():
        unwrapped = _unwrap_optional(tp)  # strip Optional[X]
        origin = get_origin(unwrapped)
        # Map: str -> pl.String, int -> pl.Int64, List[str] -> pl.List(pl.String) ...
    return schema
\end{lstlisting}

Hàm \texttt{silver\_schema\_to\_polars()} dùng \texttt{get\_type\_hints()} để
duyệt qua từng field, unwrap \texttt{Optional[X]}, và ánh xạ Python type sang
Polars dtype. \textbf{Không bao giờ} chỉnh schema bằng tay --- luôn sửa dataclass.

\subsubsection{Reader: đọc Bronze và Silver từ S3}
\label{subsubsec:silver-reader}

Hàm \texttt{get\_jobs\_silver\_by\_site()} trong \texttt{reader.py} đọc Silver
Parquet với cơ chế \textbf{chọn file mới nhất per day} bằng \texttt{LastModified}:

\begin{lstlisting}[caption={\texttt{src/storage\_layer/MinIO\_S3/layer/silver/utils/reader.py} --- đọc Silver, chỉ lấy file mới nhất per day}, label={lst:silver-reader}]
def get_jobs_silver_by_site(site, entity_name, from_date, to_date):
    """Re-runs append new parquet, khong overwrite. Moi day chi giu file
    LastModified moi nhat de tranh doc snapshot cu."""
    s3_client = get_s3_client()
    valid_s3_paths = []
    current_date = start
    while current_date <= end:
        bucket_paths = SilverBucketPaths(site, entity_name, year, month, day)
        prefix = bucket_paths.get_prefix()

        response = s3_client.list_objects_v2(
            Bucket=bucket_paths.silver_bucket_name, Prefix=prefix
        )
        parquet_objs = [
            obj for obj in response.get("Contents", [])
            if obj["Key"].endswith(".parquet")
        ]
        if parquet_objs:
            # Chi lay file moi nhat (LastModified) trong partition ngay do
            latest = max(parquet_objs, key=lambda obj: obj["LastModified"])
            latest_path = f"s3://{bucket_paths.silver_bucket_name}/{latest['Key']}"
            valid_s3_paths.append(latest_path)
        current_date += timedelta(days=1)

    # scan_parquet voi hive_partitioning=True -> source_site thanh cot that
    df_lazy = pl.scan_parquet(
        valid_s3_paths,
        storage_options=storage_options,
        hive_partitioning=True,
        missing_columns="insert",
        extra_columns="ignore",
    )
    return df_lazy
\end{lstlisting}

Đối với Bronze, \texttt{get\_jobs\_data\_from\_bronze()} dùng
\texttt{pl.scan\_ndjson()} để đọc \texttt{.jsonl.gz} trực tiếp từ S3, không cần
giải nén cục bộ.

\subsubsection{Loader: upload Parquet lên S3}
\label{subsubsec:silver-loader}

\texttt{upload\_silver\_parquet()} ghi DataFrame thành Parquet trong buffer
\texttt{BytesIO} rồi \texttt{put\_object} lên S3 --- không tạo file tạm:

\begin{lstlisting}[caption={\texttt{src/storage\_layer/MinIO\_S3/layer/silver/utils/loader.py}}, label={lst:silver-loader}]
def upload_silver_parquet(df, entity_name, source_site, date_str):
    dt = datetime.strptime(date_str, "%Y-%m-%d")
    timestamp = datetime.now().strftime("%Y%m%d_%H%M%S")

    silver_paths = SilverBucketPaths(
        source_site=source_site, entity_name=entity_name,
        year=dt.strftime("%Y"), month=dt.strftime("%m"), day=dt.strftime("%d"),
    )
    s3_key = silver_paths.get_file_key(timestamp)
    bucket = silver_paths.silver_bucket_name

    # Write Parquet to in-memory buffer (khong can file tam)
    buffer = io.BytesIO()
    df.write_parquet(buffer)
    buffer.seek(0)

    s3 = get_s3_client()
    s3.put_object(Bucket=bucket, Key=s3_key, Body=buffer.getvalue())
    return s3_key
\end{lstlisting}

\subsubsection{Cleaning pipeline}
\label{subsubsec:silver-pipeline}

Mỗi site (topcv, itviec, vietnamworks) có hàm \texttt{clean\_<site>\_jobs(df)}
riêng, áp dụng các cleaner theo thứ tự. Hàm \texttt{run\_pipeline()} trong
\texttt{pipeline.py} là shared runner:

\begin{lstlisting}[caption={\texttt{src/storage\_layer/MinIO\_S3/layer/silver/cleaning/common/pipeline.py} --- luồng chính (trích đoạn)}, label={lst:silver-pipeline}]
SILVER_SCHEMA = silver_schema_to_polars()  # compute 1 lan tai import

def run_pipeline(site, clean_fn, args):
    """Read Bronze day-by-day -> clean -> upload to Silver."""
    current_dt = start_dt
    while current_dt <= end_dt:
        date_str = current_dt.strftime("%Y-%m-%d")

        # 1. Doc Bronze cho ngay do
        lazy = get_jobs_data_from_bronze(site, args.entity_name, date_str, date_str)
        if lazy is None:
            current_dt += timedelta(days=1)
            continue

        # 2. Collect + drop rows thieu field bat buoc
        df = lazy.collect()
        df = filter_essential_rows(df)  # drop null job_title/company_name

        # 3. Apply site-specific cleaners
        cleaned = clean_fn(df)

        # 4. Enforce schema (cast/add/drop theo SilverJobItem)
        cleaned = enforce_silver_schema(cleaned, SILVER_SCHEMA)

        # 5. Upload Parquet len Silver S3 (tru khi --no_save)
        if not args.no_save:
            s3_key = upload_silver_parquet(cleaned, args.entity_name, site, date_str)

        current_dt += timedelta(days=1)
\end{lstlisting}

\texttt{enforce\_silver\_schema()} ép DataFrame khớp schema \texttt{SilverJobItem}:
cast cột tồn tại, thêm cột thiếu (null default), xoá cột không khai báo, trả
về đúng thứ tự cột. String$\rightarrow$Boolean đi qua \texttt{Int64} trung gian
(vì Polars không cast trực tiếp).

Hình~\ref{fig:s3-silver-bucket} và \ref{fig:s3-silver-drill} chụp màn hình S3 Console cho thấy Silver bucket thật với cấu
trúc \texttt{source\_site=} partition và file \texttt{.parquet}.

\begin{figure}[htbp]
    \centering
    % CHUP MAN HINH S3 CONSOLE: Bucket thethien-lakehouse-lite-silver
    % Hien thi jobs/ -> source_site=itviec/, source_site=topcv/, ...
    \includegraphics[width=\textwidth]{module2_hinh2_5.png}
    \caption{S3 Console --- Silver bucket \texttt{thethien-lakehouse-lite-silver}
    với cấu trúc Hive partition \texttt{source\_site=itviec/},
    \texttt{source\_site=topcv/}, \texttt{source\_source=vietnamworks/}.}
    \label{fig:s3-silver-bucket}
\end{figure}

\begin{figure}[htbp]
    \centering
    % CHUP MAN HINH S3 CONSOLE: Silver bucket drill-down
    % Vao jobs/source_site=itviec/year=2026/month=06/day=21/
    % Hien thi cac file clean_bronze_*.parquet
    \includegraphics[width=\textwidth]{module2_hinh2_6.png}
    \caption{S3 Console --- Silver bucket drill-down theo partition
    \texttt{source\_site=itviec/year=2026/month=06/day=21/}, hiển thị các file
    \texttt{clean\_bronze\_*.parquet}. Re-run tạo file mới (append), reader chỉ
    lấy file \texttt{LastModified} mới nhất.}
    \label{fig:s3-silver-drill}
\end{figure}


\subsection{Gold Layer: MotherDuck và DuckDB}
\label{subsec:gold-layer}

Gold Layer dùng \textbf{MotherDuck} --- dịch vụ DuckDB trên cloud. DuckDB đọc
Silver Parquet trực tiếp từ S3 qua \texttt{read\_parquet()} với
\texttt{hive\_partitioning=true}, không cần layer Polars/boto3 trung gian.

\subsubsection{MotherDuck client}
\label{subsubsec:md-client}

\begin{lstlisting}[caption={\texttt{src/storage\_layer/MotherDuck/client.py}}, label={lst:md-client}]
class MotherDuckClient:
    def __init__(self, token=MOTHERDUCK_TOKEN):
        # Ket noi MotherDuck cloud
        self.con = duckdb.connect(f'md:?motherduck_token={token}')

    def setup_s3_credentials(self):
        """Tao Persistent Secret tren MotherDuck de co quyen truy cap S3.
        Chi can chay 1 lan, MotherDuck luu an toan thong tin tren Cloud."""
        query = f"""
        CREATE OR REPLACE SECRET aws_s3_secret IN MOTHERDUCK (
            TYPE S3,
            KEY_ID '{AWS_ACCESS_KEY_ID}',
            SECRET '{AWS_SECRET_ACCESS_KEY}',
            REGION '{AWS_REGION}'
        );
        """
        self.con.sql(query)
\end{lstlisting}

\subsubsection{Đọc Silver Parquet bằng DuckDB}
\label{subsubsec:gold-read-silver}

Một single wildcard scan phủ mọi site và ngày:

\begin{lstlisting}[caption={\texttt{src/storage\_layer/MotherDuck/config.py} + \texttt{gold\_sql\_expressions.py} --- glob và SQL đọc Silver}, label={lst:gold-glob}]
# config.py
SILVER_PARQUET_GLOB = (
    f"s3://{SILVER_BUCKET}/jobs/"
    "source_site=*/year=*/month=*/day=*/*.parquet"
)

# gold_sql_expressions.py
def read_silver_parquet_sql() -> str:
    return f"""
    read_parquet(
        '{SILVER_PARQUET_GLOB}',
        hive_partitioning = true,
        filename = true,
        union_by_name = true
    )
    """
\end{lstlisting}

\subsubsection{Star schema}
\label{subsubsec:star-schema}

Gold Layer xây dựng star schema với fact table \texttt{gold.jobs} và các
child/dimension table:

\begin{longtable}{lll}
\toprule
\textbf{Table} & \textbf{Loại} & \textbf{Nguồn} \\
\midrule
\texttt{gold.jobs} & Fact & Silver scalar columns + \texttt{job\_id} + \texttt{date\_key} \\
\texttt{gold.dim\_date} & Dimension & \texttt{generate\_series()} 2023-01-01 đến 2026-05-01 \\
\texttt{gold.job\_industries} & Child & Unnest \texttt{job\_industry\_clean} (List[str]) \\
\texttt{gold.job\_benefits} & Child & Unnest \texttt{benefits\_categories\_vi} (List[str]) \\
\texttt{gold.job\_requirements} & Child & Unnest tất cả \texttt{require\_*} (List[str]) \\
\bottomrule
\caption{Gold star schema --- fact table \texttt{gold.jobs} + 4 child/dim table}
\label{tab:gold-schema}
\end{longtable}

Luồng \texttt{load\_silver\_to\_gold.py}: setup S3 credentials $\rightarrow$
ensure database/schema $\rightarrow$ build staging (dedup by
\texttt{unique\_url}, gán \texttt{job\_id}) $\rightarrow$ build fact + dims
$\rightarrow$ drop staging $\rightarrow$ log row counts.

Hình~\ref{fig:md-gold} chụp màn hình MotherDuck UI cho thấy Gold database thật với các table
và số dòng.

\begin{figure}[htbp]
    \centering
    % CHUP MAN HINH MotherDuck UI: Database lakehouse-lite, schema gold
    % Hien thi tables: jobs, dim_date, job_industries, job_benefits, job_requirements
    % Kem so dong (row count)
    \includegraphics[width=\textwidth]{module2_hinh2_7.png}
    \caption{MotherDuck UI --- database \texttt{lakehouse-lite}, schema
    \texttt{gold} với fact table \texttt{jobs} và các child/dimension table.
    Dữ liệu được DuckDB đọc trực tiếp từ Silver S3 Parquet, không qua layer
    trung gian.}
    \label{fig:md-gold}
\end{figure}


\subsection{Serving Layer: Supabase}
\label{subsec:supabase}

Supabase (Postgres) là lớp serving cho frontend. Script
\texttt{load\_silver\_to\_supabase.py} đọc Silver Parquet, ánh xạ cột cleaned
về tên gốc, và UPSERT theo \texttt{unique\_url}:

\begin{lstlisting}[caption={\texttt{src/storage\_layer/Supabase/scripts/load\_silver\_to\_supabase.py} --- ánh xạ cột và UPSERT (trích đoạn)}, label={lst:supabase-load}]
def _load_site(conn, site, from_date, to_date):
    lazy_df = get_jobs_silver_by_site(site, SILVER_ENTITY_NAME, from_date, to_date)

    # Map cleaned columns -> target names (frontend compatibility)
    _CLEANED_TO_TARGET = {
        "clean_job_title": "job_title",
        "clean_location": "location",
        "clean_company_name": "company_name",
    }

    select_exprs = []
    for col in JOB_DATA_COLUMNS:
        cleaned_src = next(
            (c for c, t in _CLEANED_TO_TARGET.items() if t == col), None
        )
        if cleaned_src:
            select_exprs.append(pl.col(cleaned_src).alias(col))
        elif col == CONFLICT_KEY:
            select_exprs.append(_unique_url_select_expr(available_columns, site))
        else:
            select_exprs.append(pl.col(col))

    df = (lazy_df.select(select_exprs)
          .filter(pl.col(CONFLICT_KEY).is_not_null()
                  & (pl.col(CONFLICT_KEY).str.strip_chars() != ""))
          .unique(subset=[CONFLICT_KEY], keep="any")
          .collect())

    # Stream UPSERT theo BATCH_SIZE chunks
    with conn.cursor() as cursor:
        for chunk in df.iter_slices(BATCH_SIZE):
            execute_values(cursor, UPSERT_SQL, chunk.rows(), page_size=BATCH_SIZE)
    conn.commit()  # per-site commit: partial progress on failure
\end{lstlisting}

Điểm đáng chú ý: ánh xạ \texttt{clean\_job\_title $\rightarrow$ job\_title}
giúp frontend nhận giá trị đã làm sạch mà \emph{không đổi} tên cột trong DB ---
schema Supabase giữ nguyên tên gốc. Mỗi site commit riêng, nên nếu một site lỗi,
các site khác vẫn đã commit.


\subsection{Vector Database: Qdrant}
\label{subsec:qdrant}

Qdrant là vector database cloud (free tier) dùng cho tính năng gợi ý việc làm
theo CV (CV-based Job Recommendation). Khác với S3/MotherDuck/Supabase lưu dữ
liệu dạng bảng, Qdrant lưu \textbf{vector embedding} + payload để tìm kiếm ngữ
nghĩa (semantic search) theo độ tương tự cosine.

\subsubsection{Mục tiêu và kiến trúc}
\label{subsubsec:qdrant-overview}

Người dùng upload CV (PDF/DOCX) $\rightarrow$ backend trích text $\rightarrow$
embed bằng Gemini $\rightarrow$ query Qdrant lấy Top~10 việc làm phù hợp. Pipeline
indexing chạy trong Lakehouse-Lite (đọc Silver Parquet $\rightarrow$ embed
$\rightarrow$ upsert Qdrant), pipeline query chạy trong JobSearchWeb (embed CV
$\rightarrow$ query Qdrant $\rightarrow$ re-rank hybrid).

\subsubsection{Thiết kế collection}
\label{subsubsec:qdrant-collection}

Collection \texttt{jobs\_v1} với vector 768 chiều, distance Cosine. Mỗi point
có ID = UUID5 từ \texttt{job\_url} (idempotent upsert), payload chứa metadata
job + mảng skills \texttt{require\_*} để filter và giải thích kết quả:

\begin{lstlisting}[caption={Thiết kế Qdrant collection \texttt{jobs\_v1}}, label={lst:qdrant-collection}]
Collection: jobs_v1
  vectors:
    size: 768
    distance: Cosine
  point id: UUID5(namespace=DNS, name=job_url)  # idempotent upsert
  payload:
    job_url            : keyword
    job_title          : text
    company_name       : keyword
    clean_location     : keyword        # filter location
    is_vietnam         : bool
    min_exp_level      : float          # filter experience
    max_exp_level      : float
    min_monthly_salary : float
    max_monthly_salary : float
    source_site        : keyword
    deadline_ts        : integer        # epoch giay; loc job het han
    require_programming_languages : keyword[]
    require_frameworks            : keyword[]
    require_tools                 : keyword[]
    require_cloud_skills          : keyword[]
    require_knowledge             : keyword[]
    require_domain_knowledge      : keyword[]
    require_foreign_languages     : keyword[]
\end{lstlisting}

Tạo \textbf{payload index} cho \texttt{clean\_location},
\texttt{min\_exp\_level}, \texttt{deadline\_ts} để filter nhanh tại query-time.

\subsubsection{Shared contract giữa 2 repo}
\label{subsubsec:qdrant-contract}

Các hằng số phải khớp tuyệt đối giữa Lakehouse-Lite (index) và JobSearchWeb
(query). Sai một cái là kết quả vô nghĩa:

\begin{longtable}{ll}
\toprule
\textbf{Hằng số} & \textbf{Giá trị} \\
\midrule
Embedding model & \texttt{gemini-embedding-001} \\
Vector dimension & 768 (MRL \texttt{output\_dimensionality=768}, L2-normalize) \\
Distance metric & Cosine \\
Qdrant collection & \texttt{jobs\_v1} \\
Task type (embed JOB) & \texttt{RETRIEVAL\_DOCUMENT} \\
Task type (embed CV) & \texttt{RETRIEVAL\_QUERY} \\
\bottomrule
\caption{Shared contract giữa Lakehouse-Lite và JobSearchWeb}
\label{tab:qdrant-contract}
\end{longtable}

Nên đặt các hằng số này vào biến môi trường dùng chung
(\texttt{EMBEDDING\_MODEL}, \texttt{EMBEDDING\_DIM},
\texttt{QDRANT\_COLLECTION}) để cả 2 repo đọc cùng một nguồn.

\subsubsection{Indexing pipeline}
\label{subsubsec:qdrant-index}

Module \texttt{src/storage\_layer/Qdrant/} đọc Silver Parquet, embed bằng
Gemini, upsert vào Qdrant. Luồng chính trong
\texttt{index\_silver\_to\_qdrant.py}:

\begin{lstlisting}[caption={\texttt{src/storage\_layer/Qdrant/scripts/index\_silver\_to\_qdrant.py} --- luồng indexing (trích đoạn)}, label={lst:qdrant-index}]
# 1. Tao collection neu chua co (idempotent), tao payload index
client.recreate_collection(
    collection_name="jobs_v1",
    vectors_config=VectorParams(size=768, distance=Distance.COSINE),
)
client.create_payload_index("jobs_v1", "clean_location",
                            PayloadSchemaType.KEYWORD)
client.create_payload_index("jobs_v1", "min_exp_level",
                            PayloadSchemaType.FLOAT)
client.create_payload_index("jobs_v1", "deadline_ts",
                            PayloadSchemaType.INTEGER)

# 2. Lap theo SITES, doc Silver bang get_jobs_silver_by_site()
for site in SITES:
    lazy = get_jobs_silver_by_site(site, SILVER_ENTITY_NAME,
                                    from_date, to_date)
    df = lazy.collect()

    # 3. Ghep text embed: title + requirements + description
    texts = [
        f"{r['clean_job_title']}\n{r['requirements_cleaned']}"
        f"\n{r['job_description_cleaned']}"
        for r in df.iter_rows(named=True)
    ]

    # 4. Embed theo batch (Gemini, RETRIEVAL_DOCUMENT, dim 768, L2-normalize)
    vectors = embed_documents(texts)  # batch + backoff khi 429

    # 5. Dung PointStruct(id=uuid5(job_url), vector=v, payload=...)
    points = [
        PointStruct(
            id=str(uuid5(NAMESPACE_DNS, r["job_url"])),
            vector=v,
            payload=asdict(QdrantJobPayload(**r)),
        )
        for r, v in zip(df.iter_rows(named=True), vectors)
    ]

    # 6. Upsert theo lo (100 diem/lan)
    client.upsert(collection_name="jobs_v1",
                  points=points, batch_size=100)

# 7. Xoa job het han: deadline_ts < now
client.delete(collection_name="jobs_v1",
              filter=Filter(must=[FieldCondition(
                  key="deadline_ts",
                  range=Range(lt=int(time.time())))]))
\end{lstlisting}

\subsubsection{QdrantJobPayload dataclass}
\label{subsubsec:qdrant-payload}

Payload được định nghĩa trong dataclass \texttt{QdrantJobPayload} --- single
source of truth của payload, tương tự \texttt{SilverJobItem}:

\begin{lstlisting}[caption={\texttt{src/storage\_layer/Qdrant/schema/data\_class.py} --- QdrantJobPayload (trích đoạn)}, label={lst:qdrant-payload}]
@dataclass
class QdrantJobPayload:
    job_url: str
    job_title: str
    company_name: str
    clean_location: Optional[str] = None
    is_vietnam: Optional[bool] = None
    min_exp_level: Optional[float] = None
    max_exp_level: Optional[float] = None
    min_monthly_salary: Optional[float] = None
    max_monthly_salary: Optional[float] = None
    source_site: str = ""
    deadline_ts: Optional[int] = None
    require_programming_languages: List[str] = field(default_factory=list)
    require_frameworks: List[str] = field(default_factory=list)
    require_tools: List[str] = field(default_factory=list)
    require_cloud_skills: List[str] = field(default_factory=list)
    require_knowledge: List[str] = field(default_factory=list)
    require_domain_knowledge: List[str] = field(default_factory=list)
    require_foreign_languages: List[str] = field(default_factory=list)
\end{lstlisting}

\subsubsection{Embedding bằng Gemini}
\label{subsubsec:qdrant-embed}

Gemini \texttt{gemini-embedding-001} (free tier) hỗ trợ đa ngôn ngữ (tiếng Anh
+ tiếng Việt), batch nhiều input/lần. Cần xử lý rate-limit (429) bằng
exponential backoff:

\begin{lstlisting}[caption={\texttt{src/storage\_layer/Qdrant/scripts/embed.py} --- embed batch với backoff (trích đoạn)}, label={lst:qdrant-embed}]
def embed_documents(texts: list[str]) -> list[list[float]]:
    """Embed batch texts voi Gemini, retry khi 429."""
    all_vectors = []
    for batch in chunked(texts, BATCH_SIZE):
        for attempt in range(MAX_RETRIES):
            try:
                resp = client.models.embed_content(
                    model="gemini-embedding-001",
                    contents=batch,
                    config=EmbedContentConfig(
                        task_type="RETRIEVAL_DOCUMENT",
                        output_dimensionality=768,
                    ),
                )
                vectors = [e.values for e in resp.embeddings]
                # L2-normalize sau khi cat dimension
                vectors = [normalize(v) for v in vectors]
                all_vectors.extend(vectors)
                break
            except RateLimitError:
                time.sleep(2 ** attempt)  # exponential backoff
    return all_vectors
\end{lstlisting}

\subsubsection{Airflow DAG}
\label{subsubsec:qdrant-dag}

DAG \texttt{index\_qdrant.py} chạy hằng ngày, dùng DockerOperator chạy
\texttt{python -m src.storage\_layer.Qdrant.scripts.index\_silver\_to\_qdrant}
trong container \texttt{lakehouse-crawler} --- đúng pattern kiến trúc, DAG chỉ
orchestrate. Lịch đặt sau khi Silver hoàn tất (nối downstream của DAG silver).

Hình~\ref{fig:qdrant-ui} chụp màn hình Qdrant UI cho thấy collection
\texttt{jobs\_v1} thật với số point và cấu hình vector.

\begin{figure}[htbp]
    \centering
    % CHUP MAN HINH Qdrant Cloud UI: Collection jobs_v1
    % Hien thi so diem (points), vector config, payload indexes
    \includegraphics[width=\textwidth]{module2_hinh2_8.png}
    \caption{Qdrant Cloud UI --- collection \texttt{jobs\_v1} với vector 768-dim
    Cosine, payload indexes cho \texttt{clean\_location},
    \texttt{min\_exp\_level}, \texttt{deadline\_ts}. Mỗi point là một job được
    embed từ Silver Parquet.}
    \label{fig:qdrant-ui}
\end{figure}


\subsection{Lưu ý quan trọng}
\label{subsec:storage-gotchas}

\begin{longtable}{p{0.42\textwidth} p{0.50\textwidth}}
\toprule
\textbf{Đặc điểm} & \textbf{Giải thích} \\
\midrule
Thư mục \texttt{MinIO\_S3} là tên legacy & Thực tế dùng boto3 $\rightarrow$ AWS S3 thật. \texttt{get\_s3\_client()} là factory duy nhất, không có \texttt{endpoint\_url}. \\
\addlinespace
Bronze upload là destructive & Nén gzip $\rightarrow$ upload $\rightarrow$ xoá cả \texttt{.gz} và \texttt{.jsonl}. Dùng \texttt{-{}-source} để cô lập một site. \\
\addlinespace
Bronze \texttt{source/} vs Silver \texttt{source\_site=} & Sự khác biệt có chủ đích: Silver cần Hive partition để DuckDB/Polars nhận \texttt{source\_site} thành cột. \\
\addlinespace
Silver re-run = append, không overwrite & Mỗi lần chạy tạo \texttt{clean\_bronze\_<timestamp>.parquet} mới. Reader chỉ lấy file \texttt{LastModified} mới nhất per day. \\
\addlinespace
Schema Silver = dataclass & Thêm/xoá field ở \texttt{SilverJobItem}, \texttt{silver\_schema\_to\_polars()} tự sinh schema. Không edit schema bằng tay. \\
\addlinespace
Gold đọc S3 trực tiếp & DuckDB \texttt{read\_parquet(hive\_partitioning=true)} --- không qua Polars/boto3. MotherDuck lưu S3 credential qua \texttt{CREATE SECRET}. \\
\addlinespace
Supabase map cleaned $\rightarrow$ raw & \texttt{clean\_job\_title $\rightarrow$ job\_title}: frontend nhận giá trị sạch, DB giữ tên gốc. UPSERT trên \texttt{unique\_url}. \\
\addlinespace
Qdrant UUID5 theo \texttt{job\_url} & Idempotent upsert: chạy index 2 lần không tăng point trùng. UUID5(namespace=DNS, name=job\_url). \\
\addlinespace
Shared contract 2 repo & \texttt{gemini-embedding-001} + 768-dim + Cosine phải khớp giữa Lakehouse-Lite (index) và JobSearchWeb (query). Đặt vào env dùng chung. \\
\addlinespace
Qdrant free tier rate-limit & Gemini 429 khi batch lớn. Gom batch + exponential backoff. Vài nghìn job/ngày nằm trong free tier. \\
\addlinespace
Bucket name global unique & Prefix \texttt{thethien-lakehouse-lite-} trong \texttt{bucket.yml}. Đổi khi share AWS account. \\
\addlinespace
Region \texttt{ap-southeast-1} & Mặc định trong \texttt{key.py}, override qua \texttt{AWS\_REGION} trong \texttt{.env}. Phải khớp region bucket. \\
\bottomrule
\caption{Lưu ý quan trọng của Storage Layer}
\label{tab:storage-gotchas}
\end{longtable}


\subsection{Lệnh chạy (verification)}
\label{subsec:storage-commands}

\begin{lstlisting}[caption={Lệnh chạy từng stage (từ repo root)}, label={lst:storage-commands}]
# 1. Bronze: gzip + upload temp JSONL -> S3
python -m src.storage_layer.MinIO_S3.layer.bronze.main --source topcv

# 2. Silver: clean Bronze -> Parquet (dry run, khong ghi S3)
python -m src.storage_layer.MinIO_S3.layer.silver.cleaning.clean_topcv.main_process \
    --from_date 2026-06-21 --to_date 2026-06-21 --no_save

# 3. Silver: clean + upload Parquet
python -m src.storage_layer.MinIO_S3.layer.silver.cleaning.clean_topcv.main_process \
    --from_date 2026-06-21 --to_date 2026-06-21

# 4. Supabase: UPSERT Silver -> Postgres
python -m src.storage_layer.Supabase.scripts.load_silver_to_supabase \
    --from_date 2026-06-21 --to_date 2026-06-21

# 5. Gold: DuckDB read S3 Parquet -> MotherDuck star schema
python -m src.storage_layer.MotherDuck.scripts.load_silver_to_gold

# 6. Qdrant: embed Silver -> vector DB (backfill hoac daily)
python -m src.storage_layer.Qdrant.scripts.index_silver_to_qdrant \
    --from_date 2026-06-21 --to_date 2026-06-21
\end{lstlisting}

Cờ \texttt{-{}-no\_save} (Silver) cho phép dry-run: chạy toàn bộ pipeline làm
sạch mà không ghi lên S3, hữu ích để debug. Cờ \texttt{-{}-export\_parquet} dump
Parquet cục bộ để kiểm tra bằng tay.
